\chapter{Dentures}

\section*{Overview}
Dentures are removable replacements for missing teeth. A complete denture replaces all teeth in the upper or lower jaw; a partial denture replaces some teeth and may attach to remaining teeth with clasps or other supports. Dentures can help with chewing, speech, and appearance. They rest on the gums or, in some designs, attach to dental implants.

\section*{Alternative Names}
False teeth; removable dental prosthesis; complete or full denture; partial denture.

\section*{Principle}
A dentist or dental prosthetist makes a denture to fit the shape of the mouth and replace missing teeth. A conventional denture is made after the gums have healed following tooth removal. An immediate denture is made in advance and fitted on the day teeth are removed, so the person is not without teeth while healing. Because the gums and jaw change shape as they heal, an immediate denture often needs adjustment, relining, or replacement later.

\section*{Why It Is Done}
A denture may replace some or all missing teeth when a person wants to restore chewing or speech, or prefers a removable option. Bridges and implants are other options. A dental professional discusses which choices fit the person's mouth, health, needs, and preferences.

\section*{Role in Care}
Dentures are restorative devices, not a test or surgery. The fitting process includes examinations and measurements or impressions so the denture can be made and adjusted to the mouth.

\section*{How It Is Performed}
The dental professional examines the mouth, discusses the design, and takes impressions or digital scans and bite measurements. Trial fittings may be used to check the shape, fit, bite, and appearance before the final denture is made. At delivery, the clinician checks the fit and explains how to insert, remove, and care for it. Follow-up visits may be needed to adjust sore or tight areas. An implant-supported denture requires implant surgery and additional healing and fitting steps.

\section*{Preparation}
Tell the dental team about medical conditions, medicines, allergies, and any teeth that are painful or loose. Discuss whether teeth need to be removed first, whether an immediate denture is an option, what the denture will be made from, and the expected number of visits. If extractions are planned, ask for instructions about healing and the first days with the denture.

\section*{Contraindications and Precautions}
Most people can use dentures, but the design may need to account for sore gums, oral infection, dry mouth, jaw changes, or difficulty handling and cleaning the denture. Treating active oral disease first may improve comfort and fit. Partial dentures need regular care of the remaining teeth and gums, since plaque around the denture can contribute to decay or gum disease.

\section*{Risks}
New dentures may feel bulky and can cause temporary sore spots, extra saliva, or difficulty speaking and chewing while the person adjusts. A poor fit can cause persistent pain, mouth ulcers, or trouble eating. Dentures may loosen as the shape of the gums and jaw changes. They can crack if dropped, and partial dentures can affect the teeth that support them if they are not kept clean or checked.

\section*{Aftercare and Recovery}
It can take days or weeks to get used to new dentures. Start with soft foods, take small bites, and practise speaking; follow the clinician's advice about how long to wear them, especially after extractions. Brush the denture daily with a soft denture brush and a cleaner labeled as safe for its material. Use soaking cleansers only after removing the denture from the mouth, and follow the product directions. Clean the gums, tongue, and any remaining teeth as well. In general, remove dentures at night to let the mouth rest and store them as directed. After an extraction, the dentist may instruct a person to keep an immediate denture in place overnight or for the first 24 hours; follow the individual instructions. Avoid hot or boiling water and household bleach unless the denture's instructions specifically say they are safe.

\section*{Outcome and Follow-up}
The dental professional checks comfort, retention, and the bite and adjusts pressure areas as needed. Adhesive can help a well-fitting denture stay in place, but it is not a substitute for assessment if a denture is loose or painful. Use adhesive only as directed and remove it from the denture and mouth each day. Needing more adhesive or applying it repeatedly during the day can mean the denture needs adjustment or replacement. The gums and jaw can change over time, so a denture may need relining, repair, or replacement.

\section*{Expected Results}
A well-fitting denture can improve chewing, speech, and appearance, but it may not feel or function exactly like natural teeth. Comfort and stability vary, and some people need practice or adjustments. Continue regular dental visits even if all natural teeth have been removed, so the mouth and denture can be checked.

\section*{Follow-up}
Attend fitting and adjustment visits and routine dental examinations. Ask a dentist to check a denture that becomes loose, cracks, causes ongoing soreness, or needs increasing amounts of adhesive. Keep the denture clean and follow the product's storage instructions when it is out of the mouth.

\section*{When to Seek Medical Care}
Contact a dentist if a denture causes ongoing pain or ulcers, becomes loose or breaks, or makes it difficult to eat. Seek urgent care for facial swelling, fever with mouth swelling, or difficulty swallowing or breathing.

\section*{References}
\begin{enumerate}
  \item NHS. Dentures. https://www.nhs.uk/tests-and-treatments/dentures/
  \item American Dental Association. Denture Care and Maintenance. https://www.ada.org/resources/ada-library/oral-health-topics/dentures
  \item Guy's and St Thomas' NHS Foundation Trust. Dentures. https://www.guysandstthomas.nhs.uk/health-information/dentures
  \item Cambridge University Hospitals NHS Foundation Trust. Denture care: Looking after your dentures. https://www.cuh.nhs.uk/patient-information/denture-care-looking-after-your-dentures/
\end{enumerate}

\clearpage
